%%
% BIThesis 研究生学位论文模板 The BIThesis Template for Graduate Thesis
% This file has no copyright assigned and is placed in the Public Domain.
%%

\begin{conclusion}

本文围绕面向垂直领域知识图谱构建的高质量可控知识抽取问题，针对领域标注数据稀缺、抽取结果难以稳定入库以及本地部署成本较高等实际困难，构建了一条贯通“数据—抽取—应用”的完整技术路线，并以军事新闻领域为实例完成了方法设计、实验验证和系统实现。研究结果表明，现如今，知识抽取的关键已不再只是能否抽取的问题，而在于能否同时兼顾训练数据质量、结果结构合法性、部署可行性和下游可用性。围绕这一目标，本文首先提出了基于 SO-CoT 自优化与多模型事实评估的银标数据构建方法，通过提示自优化、事实校验与回流纠错，形成了包含 10647 份军事新闻报道和 40565 个实体关系三元组的领域银标数据集；进一步地，提出了面向轻量化本地部署的参数高效本体约束知识抽取方法，将 QLoRA 领域适配、本体约束和多候选评分统一纳入同一流程，在 1000 条测试样本上的 F1 达到 74.95\%；最后，完成了知识图谱构建与问答原型系统实现，在 1000 条输入文本上构建出 4546 个节点、6017 条关系的领域知识图谱，验证了本文方法在稳定入库、基础查询和简单问答场景中的有效性。

本文的创新性主要体现在三个方面。第一，本文将高质量银标数据构建从训练后的清洗补救前移为训练前的核心前提，通过 Teacher 端闭环机制把数据质量控制前置，缓解了垂直领域人工标注成本高、自动标注噪声大的问题。第二，本文在 Student 端将小参数模型、QLoRA、本体约束和多候选选择协同设计，使模型不仅具备领域适配能力，而且具备更好的结构合法性和稳定入库能力。第三，本文没有将研究停留在句子级抽取指标上，而是将知识图谱构建和下游问答验证纳入整体评估，使方法评价从单一精度提升扩展到工程可用性与知识服务能力，更贴近真实场景中的领域知识组织需求。

从应用前景看，本文提出的方法并不限于军事新闻场景，而是对医学、金融、科技等同样具有多源异构、术语密集和持续更新特征的垂直领域文本具有较强参考价值。其现实意义在于：一方面，能够以较低的人力成本提升领域知识资源的结构化组织效率，为专业资讯管理、知识服务和行业数据治理提供技术支撑；另一方面，能够在本地部署条件下形成可持续运行的知识抽取能力，降低长期依赖超大模型在线推理带来的算力和运维压力。就社会价值而言，该研究有助于提高专业资讯的可检索性、可追溯性和可利用性，支撑智能问答、知识检索和辅助分析等服务形态；就经济价值而言，该研究有望降低领域知识图谱构建中的数据标注、模型部署和后处理成本，提高知识资产沉淀与复用效率，具有较好的工程推广潜力。

同时也应看到，本文研究仍有进一步拓展空间。首先，当前方法主要在军事新闻场景下完成验证，后续仍需在更多领域和更多文本风格下检验其跨领域迁移能力与鲁棒性。其次，本文本体设计和抽取目标目前以实体关系三元组为主，未来可进一步面向事件级知识表示、跨句关系、时序演化和增量更新展开研究，以增强知识图谱对复杂事实链条的表达能力。再次，本文虽然通过本体约束和多候选评分提升了结果可控性，但在长文本、难例样本和高频动态更新场景下，抽取效率与自动纠错能力仍有优化空间，后续可继续加强人机协同反馈、结构约束与持续学习机制之间的结合。总体而言，本文完成了从高质量数据构建、可控抽取到图谱应用验证的一体化探索，为面向新闻领域知识图谱构建的知识抽取研究提供了一条较完整、可落地的实现路径，也为后续进一步走向更广领域、更复杂知识形态和更强工程实用性奠定了基础。

\end{conclusion}
